\chapter{Teori og beregningsgrunnlag}
\label{ch:theory}
\section{Sanntid, tasking og exceptions i Ada}
Et sanntidssystem vurderes både etter resultatet og tidspunktet resultatet
foreligger. En deadline må begrunnes i prosessen som styres. For bilen
avhenger en forsvarlig reaksjonsfrist av fart, avstand og bremseevne.
En innstilt ventetid eller et vellykket bygg er ikke et deadlinebevis.

Ada gir avgrensede typer, pakker, exceptions og eksplisitt samtidighet.
Et protected object kan beskytte kortvarige operasjoner på delt tilstand;
beskyttelsen erstatter ikke planlegging av lange eller ventende operasjoner.
Tasking og synkronisering beskrives i språkets referansemanual
\cite{adarm,adaprotected2026}. I løsningen publiserer motor-/kontrollkoden
lystilstand, mens én task utfører GPIO-skriving. Ingen I2C, logging eller
\kode{delay} skjer inne i protected object.

En exception som ikke håndteres lokalt, forplanter seg til en omgivende handler
innen samme task. En feil som avslutter en separat task, blir ikke en vanlig
exception i hovedprogrammet \cite{adaexceptions2026}. Lys-tasken har derfor
sin egen feilgrense og publiserer feil før den parkeres. Hovedprogrammet
oppdager statusen og stopper gjennom sin egen feilgrense. Kjøretidsstøtten
avhenger av valgt runtime \cite{adart2026}; prosjektet bruker full
embedded-runtime og aktiverer assertions og overflow-sjekker i egen kode.

\section{Ultralydmåling og usikkerhet}
HC-SR04 bruker en triggerpuls og returnerer en ECHO-puls som representerer
tur-retur-tiden til en reflektor. Avstanden modelleres som
\begin{equation}d=\frac{c\,\Delta t}{2},\end{equation}
der $c$ er lydhastigheten. Ved antatt $c=343$ m/s blir én centimeter omtrent
58,3 mikrosekunder ECHO-tid. Temperatur, reflektorvinkel, overflate,
minsteavstand og elektrisk oppsett påvirker resultatet. Databladet anbefaler
mer enn 60 ms målesyklus \cite{hcsr04}; prosjektet bruker global 65 ms
triggeravstand for de tre sensorene og forhindrer samtidig utsending.

Ingen ECHO, fast høy inngang eller tidsoverskridelse betyr ukjent avstand.
Verdien 0 får aldri representere fri bane. Prosjektet godtar bare 2--200 cm
i beslutninger. Ved stasjonær filtrering må alle tre prøver være gyldige,
minst ett par være innen 15 cm, og den minste av de tre beholdes.
Et kort ekko kastes dermed ikke fordi to andre ser lenger. Dette er et
konservativt konsistensfilter, ikke en median og ikke en sannhetsgaranti.

\section{Mecanum og motorkommandoer}
Skrå ruller gir Mecanum-hjul både langsgående og sideveis kraftbidrag.
Et korrekt montert firehjulsoppsett kan kombinere translasjon og rotasjon
\cite{mecanum}. Den aktive løsningen bruker rett frem og rotasjon til
venstre/høyre. Sidelengs kjøring og ryggemanøvrer inngår ikke i demonstrert
autonomi. Hjulplassering, rulleretning og motorfortegn må kontrolleres
fysisk; en ideell kinematisk matrise alene dokumenterer ikke dette.

DFR0548 mottar PWM-registerverdier via I2C \cite{dfr0548}. PWM 800 er en
åpen kommando, ikke en måling av fart. Batterispenning, last, underlag og
friksjon endrer bevegelsen ved samme kommando. Derfor kan verken
kjørelengde eller svingvinkel beregnes sikkert fra tidsvarighet alene.

\section{Korridor, søkekostnad og stopplengde}
To stråler med vinkelavstand $\Delta\theta$ og minste målt avstand $d$
brukes i en enkel korridormodell:
\begin{equation}
 w=2d\sin(\Delta\theta/2),\qquad
 w_{\mathrm{krav}}=16+2\cdot5=26\text{ cm}.
\end{equation}
Ved $\Delta\theta=12^\circ$ kreves $d\ge124{,}4$ cm. Beregningen gir en
konkret avvisningsregel. Den forutsetter en tilstrekkelig sammenhengende
åpning og beskriver ikke hele området mellom sparsomme ultralydstråler
eller hele bilens sveipareal. Sentrerte målinger og sideklaring kreves i
tillegg. Smale eller absorberende hindringer kan fortsatt overses.

For to uavhengige søkekandidater med kostnad $c_i$ og sannsynlighet $p_i$
for aksept er forventet kostnad ved rekkefølgen $1,2$ lik
$c_1+(1-p_1)c_2$. Sammenligning med motsatt rekkefølge gir størst
$p_i/c_i$ først. Prosjektet har ingen målte $p_i$; størst sideklaring først
er en begrunnet heuristikk, ikke dokumentert global optimalitet.

En bremsemodell med konstant retardasjon $a_b>0$ gir
\begin{equation}\label{eq:stop}
 d_{\mathrm{stopp}}=v\,t_r+\frac{v^2}{2a_b}+m.
\end{equation}
For tilgjengelig avstand $d>m$ følger
$v_{\max}=\sqrt{(a_b t_r)^2+2a_b(d-m)}-a_b t_r$.
Dette er analysegrunnlag; modellen er ikke en implementert fartsregulator.
Fart, konservativ retardasjon, reaksjonstid og margin må måles før
den kan begrunne bilens driftsområde.
